\chapter{第四讲\quad 犯罪主体}
\label{lecture:04}
\noindent{\small\hyperref[ref:lecture:04]{查阅本讲扩展规则}}\par

\section{第一节\quad 自然人}

\paragraph{定义\quad 定罪身份与量刑身份}
设 \(\mathrm{Role}_g(a,x)\) 表示行为人 \(a\) 在实施行为 \(x\) 时具备罪名 \(g\) 要求的特殊身份。对真正身份犯，若 \(a\) 是拟评价的直接实行者，身份须加入构成要件。量刑身份则作为进入特定量刑分支的附加前提。设 \(K(a,x)\) 为其余要件合取，简写为
\begin{align*}
 \text{定罪身份}\qquad \mathrm{Elem}_g(a,x)
  &:=K(a,x)\land\mathrm{Role}_g(a,x),\\
 \text{量刑身份}\qquad \mathrm{Elem}_g(a,x)
  &:=K(a,x),\\
 \mathrm{Elem}_g(a,x)\land\mathrm{Role}_g(a,x)
  &\Rightarrow \text{适用相应的法定量刑规则}.
\end{align*}
例如，第247条刑讯逼供罪以司法工作人员为主体。第243条诬告陷害罪的国家机关工作人员身份进入从重分支，第254条报复陷害罪中的国家机关工作人员身份则进入定罪要件。行为发生后取得的身份通常不能反向补足当时的定罪身份。某些角色本身由行为形成，例如组织、领导某组织，不能只因名称带有角色就一概归为真正身份犯。

\paragraph{推导\quad 角色缺项只切断相应的实行分支}
设甲为执行职务的司法工作人员，为取口供实施刑讯，满足其余要件。身份与行为各项合取支持第247条的实行分支。将甲改为独自对他人殴打取信息的普通人，身份项为假，不能由同一规则认定刑讯逼供罪，仍应另查故意伤害等罪名。又设普通人丙明知甲的职务行为及刑讯计划，故意提供帮助，教材的共犯分析先验甲的违法实行，再验丙的帮助及共同故意。缺少司法工作人员身份只能否定丙作为该罪独立直接实行者的这一条路径，不能单凭该点排除共犯。身份在不同参与方式中的作用须按具体罪名和解释规则判断。

\paragraph{国家工作人员\quad 把公务放回行为时点}
令 \(\mathrm{PublicDuty}(a,x)\) 表示涉案事务具有公务性质。来源分四类。\(S_1\) 表示国家机关人员，\(S_2\) 表示国有单位等法定机构人员，\(S_3\) 表示由法定机构向非国有单位等委派，\(S_4\) 表示其他依法从事公务的来源。第93条的筛选可概括为
\[
 \mathrm{State}(a,x)
  :=\mathrm{PublicDuty}(a,x)\land(S_1\lor S_2\lor S_3\lor S_4),
\]
并以条文各分支的具体主体范围为准。公务着重考察职责来源、公共管理和监督管理公共事务、国有财产等内容。正式编制与劳动合同身份可作为证据线索，结论须落在涉案职权和实际职责。村基层组织人员协助政府管理征地补偿费用时，可以按第93条第2款立法解释进入公务分支，纯粹办理本村自治事务须另作判断。贪污罪还有第382条第2款规定的受委托管理、经营国有财产人员这一入口，不能把全部贪污主体压缩成第93条一项。

\paragraph{同人两事的对照推导}
自拟题设。甲在公立医院负责公共采购审批，利用审批权为供应商谋利并违法收财。甲另一次仅以医生身份开处方，利用开处方便利为销售方谋利并收财。两次均设证据充分、各自达到相应入罪条件。根据第93、163、385条及两高关于办理商业贿赂刑事案件的意见，第一次的公共采购管理与职权来源支持 \(\mathrm{State}(\text{甲},x_1)\)，再联结职务便利与收财谋利事实，审查受贿罪。第二次题设仅给技术性处方权，无法套用第一次的公务前提，依非国家工作人员受贿分支审查。若第二次还利用行政管理权影响公立医院采购，需重审职务关联。保持人员姓名相同并不能保证身份评价在两事件中相同。

\section{第二节\quad 单位犯罪}

\paragraph{定义\quad 先问法条是否开放单位入口}
令 \(\mathrm{UnitCap}(u)\) 表示组织 \(u\) 可作为刑法单位主体，\(\mathrm{AllowsUnit}(g)\) 表示分则或其他法律允许该罪追究单位，\(\mathrm{Attr}(u,x)\) 表示行为可归属于单位，\(\mathrm{Elements}_g(u,x)\) 表示该单位分支的其余要件成立。第30条和教材的归属审查可组合为
\[
 \mathrm{UnitOff}(g,u,x)
 \;:=\;\mathrm{UnitCap}(u)\land\mathrm{AllowsUnit}(g)
      \land\mathrm{Attr}(u,x)\land\mathrm{Elements}_g(u,x).
\]
最后一项包含该罪要求的主客观要件及门槛。单位意思由何人、依何种权限形成，行为以谁名义实施，利益怎样归属，均是审查 \(\mathrm{Attr}\) 的材料。教材用「单位意志／个人意志」与「单位利益／个人利益」的两轴区分典型案例。利益流向用于判断归属，不能机械设定每个单位罪都要求本单位获利。特定罪还可能有单位过失分支，需依据其法条认定。

单位可犯、个人不可犯的典型是单位行贿罪，单位与自然人均可进入的典型是生产、销售伪劣产品罪，单位入口未开放的典型是普通盗窃罪。对主体资格，法释〔1999〕14号第1条按组织类型列举，具有法人资格是其中独资、私营等组织分支的条件。教材对分支机构或内设机构另查是否以该机构名义实施、所得归该机构，以及所涉罪的具体规定。法人登记、机构名称和单位归属是不同问题。

\paragraph{推导\quad 单位决策与成员擅自行为}
某依法设立的公司由有权机关决议，以公司名义批量生产、销售伪劣产品，收益进入公司，销售额及其余要件满足单位分支。第150条提供第140条所涉犯罪的单位入口，故 \(\mathrm{AllowsUnit}=1\)。授权、名义和归益共同支持 \(\mathrm{Attr}=1\)。四项成立，进入单位责任。将事实改为职员盗用公司名义、所得由实施者私分，法释〔1999〕14号第3条指向自然人犯罪，单位归属失败。再看该解释第2条。个人为从事违法犯罪而设立机构并实施犯罪，或者机构成立后以犯罪为主要活动，均触发排除单位犯罪的规则。以正常经营中偶发违法为题设时，须审查本次行为归属，不能倒推其设立即为犯罪。

\paragraph{反例\quad 缺少单位入口与单位门槛不足}
在甲案，领导集体决定盗电。依第30条和2014年全国人大常委会对第30条的解释，盗窃没有单位追责入口，须分别评价组织、策划、实施者个人。对每个人仍要核对行为、故意、参与方式及法定门槛。单有职务或列席会议无法推出个人有罪。

在乙案，题设规定某罪存在单位和自然人两个入口，且对同一计量指标分别设单位门槛 \(T_u\) 和个人门槛 \(T_p\)。若该案已认定归属于单位，且 \(T_p\le m<T_u\)，则单位分支的数额要件失败。个人门槛较低这一事实无法自行证明该行为属于某自然人的独立个人犯罪。模型使用抽象门槛，具体标准应查对应罪名的现行规范。教材据此区分「法律没有单位入口」与「有入口而本案未达该分支门槛」。单位与单位外主体是否构成共同犯罪，要另查各自入罪入口、共同故意与行为。仅因多名员工执行单位决策，无法把单位本身与内部人员自动认定为两个独立的共犯主体。

\paragraph{处罚\quad 定罪层与执行主体分开}
依第31条，原则上对单位判处罚金，对直接负责的主管人员和其他直接责任人员依法处刑。若分则或其他法律另有规定，按特别规定。单罚分支仍需辨认其处罚对象与依据。最高人民法院关于适用《中华人民共和国刑事诉讼法》的解释第344条区分两类诉讼事实。被吊销执照、破产但尚未完成清算和注销登记的，继续审理。单位已被撤销或注销的，对直接负责的主管人员和其他直接责任人员继续审理。单位后续状态应单列，不能由注销倒推原案从未成立单位犯罪。

\subsection*{闭卷复述与检验}
\textbf{口诀}\quad「定罪身份入要件，量刑身份入幅度。人看当时办何事，单位先看法条开门，再验归属和门槛。」

\textbf{复述路线}\quad
把自然人和单位分两列。自然人列填行为时点、所从事务、身份作用及参与方式。单位列填资格、法定入口、决策权限、名义与收益、数额或情节。把法条缺入口和本案数额不达标用两种符号标出，最后才填处罚对象。

\begin{enumerate}
\item 国家机关工作人员犯诬告陷害，与普通人犯同罪的身份区别落在哪一层？\\
\textbf{答}\quad 在第243条的量刑层。只改变身份仍可能成立基础罪，具备法定身份再进入从重分支。若改为第254条的报复陷害罪，身份进入构成。
\item 某机构依集体决议实施无单位入口的罪，全部成员是否当然按自然人犯罪处罚？\\
\textbf{答}\quad 依第30条立法解释转入个人审查。逐人认定组织、策划、实行或其他参与，以及该人的故意、责任等条件。集体决议不能代替个人要件。
\item 有单位入口的罪中，单位数额不达标但超过个人门槛。能否只换阈值定责任人员的个人犯罪？\\
\textbf{答}\quad 不能。单位分支因门槛缺项而不成立。独立个人犯罪需另有归属于个人的行为事实及其余要件，不能只替换数额比较式。
\end{enumerate}
